\begin{lemma}
For any fixed red row map $\rho$, consider the set of admissible injections $\text{Fin } n \hookrightarrow \text{Fin } m \times \text{Fin } m$ whose red-coordinate map is $\rho$. If this set is nonempty, then for any fixed set $T$ of original vertices and any fixed target set $S$ of blue coordinates, the fraction of such injections that assign all vertices in $T$ to blue coordinates in $S$ is bounded by $(|S|/m)^{|T|}$.
\end{lemma}
\begin{proof}
Let
\[
X_{\rm all}
=\{\phi:\operatorname{Fin}n\hookrightarrow
\operatorname{Fin}m\times\operatorname{Fin}m:
(\phi(v))_1=\rho(v)\text{ for every }v\}
\]
and let \(X_{\rm good}\subseteq X_{\rm all}\) be the subset for which
\((\phi(v))_2\in S\) for every \(v\in T\).  By hypothesis
\(X_{\rm all}\) is nonempty.  For each row \(i\in\operatorname{Fin}m\), put
\[
R_i=\{v:\rho(v)=i\},\qquad q_i=|R_i|,\qquad
T_i=T\cap R_i,\qquad h_i=|T_i|.
\]
Nonemptiness of \(X_{\rm all}\) implies \(q_i\le m\) for every \(i\): in an
admissible injection, the \(q_i\) vertices of \(R_i\) all have first
coordinate \(i\), so injectivity forces their second coordinates to be
distinct elements of \(\operatorname{Fin}m\).

An admissible injection is equivalently, for each row \(i\), a choice of
distinct blue coordinates for the vertices in \(R_i\).  The row choices are
independent across different rows because pairs with different first
coordinate are automatically distinct.  Thus the number of all row choices
in row \(i\) is the falling factorial
\[
(m)_{q_i}=m(m-1)\cdots(m-q_i+1),
\]
and \(|X_{\rm all}|\) is the product of these row counts.

Fix a row \(i\).  We prove that the fraction of row choices sending every
vertex of \(T_i\) into \(S\) is at most \((|S|/m)^{h_i}\).  First suppose
\(h_i>|S|\).  The \(h_i\) vertices in \(T_i\) must receive distinct blue
coordinates, since the full map is injective and all of them have the same
red coordinate \(i\).  They cannot all receive values in the set \(S\), which
has only \(|S|\) elements.  Hence the favorable row count is \(0\), the row
ratio is \(0\), and the desired row bound follows from nonnegativity of
\((|S|/m)^{h_i}\).

It remains to handle the case \(h_i\le |S|\).  First choose the blue
coordinates of the \(h_i\) constrained vertices in \(T_i\).  These must be
distinct elements of \(S\), giving \((|S|)_{h_i}\) choices.  After those
choices are made, the remaining \(q_i-h_i\) vertices in \(R_i\setminus T_i\)
may use any blue coordinates not already used, giving
\((m-h_i)_{q_i-h_i}\) choices.  Therefore the favorable row ratio is
\[
\frac{(|S|)_{h_i}(m-h_i)_{q_i-h_i}}{(m)_{q_i}}.
\]
Since \(h_i\le q_i\le m\), the denominator factors as
\[
(m)_{q_i}=(m)_{h_i}(m-h_i)_{q_i-h_i},
\]
so the row ratio is
\[
\frac{(|S|)_{h_i}}{(m)_{h_i}}
=\prod_{\ell=0}^{h_i-1}\frac{|S|-\ell}{m-\ell}.
\]
Because \(S\subseteq\operatorname{Fin}m\), we have \(|S|\le m\), and by the
current case \(h_i\le |S|\).  Applying
\(\noderef{OppositeProjectionFallingFactorialBound}\) with
\((S,m,h)=(|S|,m,h_i)\) gives
\[
\prod_{\ell=0}^{h_i-1}\frac{|S|-\ell}{m-\ell}
\le (|S|/m)^{h_i}.
\]
This proves the same row bound in the second case.

All row ratios are nonnegative, so multiplying the rowwise inequalities gives
\[
\frac{|X_{\rm good}|}{|X_{\rm all}|}
\le \prod_i (|S|/m)^{h_i}
=(|S|/m)^{\sum_i h_i}.
\]
Finally, the sets \(T_i=T\cap R_i\) are pairwise disjoint and their union is
\(T\), since every vertex \(v\in T\) has exactly one row \(\rho(v)\).  Hence
\(\sum_i h_i=|T|\), and the last display is exactly
\[
\frac{|X_{\rm good}|}{|X_{\rm all}|}
\le (|S|/m)^{|T|}.
\]
\end{proof}
